\documentclass[10pt,conference]{IEEEtran}
\usepackage[T1]{fontenc}
\usepackage{booktabs}
\usepackage{microtype}
\usepackage{url}
\usepackage[hidelinks]{hyperref}
\usepackage{balance}

\title{One Byte, One Rank Reversal: Auditing HumanEval Pipeline Sensitivity for Base Code Models}
\author{\IEEEauthorblockN{Anonymous Author(s)}}

\begin{document}
\maketitle

\begin{abstract}
Published benchmark tables are often reused as though model rankings were
properties of checkpoints alone. We test whether the HumanEval order of five
open base-model identifiers in the Qwen2.5-Coder report transfers to one fully
disclosed local pipeline. Each frozen checkpoint produced one greedy completion
for all 164 tasks in BF16 on a consumer GPU; EvalPlus 0.3.1 scored HumanEval and
HumanEval+ in a network-isolated container. The published order failed to
transfer: StarCoder2-3B scored 3/164 rather than the displayed 31.7\%, reversing
its order with Qwen2.5-Coder-0.5B. Kendall's $\tau_b$ was 0.8 and the paired
StarCoder2-minus-Qwen difference was $-21.95$ percentage points (task-resampling
interval $[-28.66,-15.85]$). A stock EvalPlus control scored 2/164. In a post-hoc
counterfactual, deleting only the trailing newline appended to stripped
base-model prompts raised StarCoder2 to 49/164 and restored the five-model
order. A completed $2\times2$ newline-by-stop factorial scored its fourth cell
at 50/164. The result shows that an apparently innocuous prompt byte can change
a base-model score by about 28 points and reverse a portable-looking rank.
\end{abstract}

\section{Introduction}
HumanEval measures whether generated Python functions satisfy hidden tests
\cite{chen2021humaneval}; HumanEval+ strengthens those tests
\cite{liu2023evalplus}. Reports frequently place several checkpoints in one
table, inviting reuse of their displayed order. Yet functional correctness
depends on prompts, decoding, stopping, library behavior, sanitization, and
execution policy as well as weights. For base models, a boundary byte can alter
whether a decoder continues a function or begins a new top-level definition.

We audit one concrete transport claim: does the HumanEval order of five base
models in Qwen2.5-Coder Table 5 \cite{hui2024qwen25coder} transfer when the same
model identifiers are run through a fully disclosed local stack? This is a
rank-transfer study, not an exact rerun: the source report does not identify
immutable checkpoint revisions or every inference detail. Our contributions are:
(1) a prespecified rank endpoint with paired task-resampling analysis; (2) a
complete, container-evaluated five-model replication package; and (3) a
post-hoc mechanism audit showing that one trailing newline is sufficient to
explain most of a 29.9-point StarCoder2 discrepancy and the observed reversal.

\section{Design}
\subsection{Target and decision rule}
The published low-to-high order is Qwen2.5-Coder-0.5B, StarCoder2-3B,
DeepSeek-Coder-1.3B, Qwen2.5-Coder-1.5B, and Qwen2.5-Coder-3B. StarCoder2's
displayed 31.7\% HumanEval and 27.4\% HumanEval+ values also occur in its own
report \cite{lozhkov2024starcoder2}; we treat these appearances as one
measurement, not independent replications.

The primary statistic is Kendall's $\tau_b$ between the published order and
local greedy pass@1. We paired models by task and drew 10,000 percentile
task-resampling replicates with seed 2026. We declared transfer only if the
order was identical ($\tau_b=1$) and no adjacent pair reversed. We declared
failure if a pair reversed and its paired 95\% interval excluded zero in the
reversed direction; other outcomes were inconclusive. The question, endpoint,
and rule were committed before aggregate correctness was inspected.
The supplement retains the specification text. This repository record is not
an independent preregistration. The four adjacent intervals are unadjusted;
they are task-composition summaries, not simultaneous confidence statements.

\subsection{Frozen pipeline and safety}
We resolved the five public identifiers to immutable model revisions. All used
the EvalPlus base-model HumanEval prompt, forced base mode, greedy decoding,
one completion per task, 512 new tokens, BF16 weights, and an NVIDIA RTX 4070.
Generation used PyTorch 2.8.0+cu128 and Transformers 4.57.6. The same completion
was scored on all 164 HumanEval tasks and HumanEval+ v0.1.10 with EvalPlus 0.3.1
at a pinned revision. Generated Python ran only inside a disposable Linux
container with networking disabled, credentials absent, capabilities dropped,
and no writable host mount beyond the result exchange directory.

Two equivalence checks guard wrapper effects. A batch-one stop criterion was
compared against the optimized decode-once implementation before the full run.
We also regenerated two tasks for every checkpoint through unmodified EvalPlus;
raw and sanitized records were byte-identical to retained primary records.

\section{Primary Results}
Table~\ref{tab:primary} gives exact counts. The four non-StarCoder positions
retained their relative order, but StarCoder2 fell below Qwen-0.5B. Kendall's
$\tau_b$ was 0.8 with task-resampling interval $[0.6,0.8]$. StarCoder2 minus
Qwen-0.5B was $-21.95$ points, interval $[-28.66,-15.85]$, so the prespecified
decision was \emph{failed to transfer}. HumanEval+ preserved the same local
order. Exact percentage equality was never the endpoint.

\begin{table}[t]
\caption{Published pass@1 (\%) and local passes out of 164. HE+: HumanEval+.}
\label{tab:primary}
\centering
\footnotesize
\setlength{\tabcolsep}{3.2pt}
\begin{tabular}{@{}lrrrr@{}}
\toprule
Model & Pub. HE & Local HE & Pub. HE+ & Local HE+ \\
\midrule
Qwen-0.5B & 28.0 & 39/164 & 23.8 & 33/164 \\
StarCoder2-3B & 31.7 & 3/164 & 27.4 & 3/164 \\
DeepSeek-1.3B & 34.8 & 56/164 & 26.8 & 47/164 \\
Qwen-1.5B & 43.9 & 64/164 & 36.6 & 56/164 \\
Qwen-3B & 52.4 & 86/164 & 42.7 & 70/164 \\
\bottomrule
\end{tabular}
\end{table}

A full StarCoder2 run through the unmodified EvalPlus generation path scored
2/164 on both suites. The optimized wrapper therefore does not explain the low
score or reversal.

As an additional, post-hoc multiplicity sensitivity check, the reversed pair
has 37 tasks passed only by Qwen-0.5B and one passed only by StarCoder2.
A two-sided exact McNemar test gives $p=2.84\times10^{-10}$; Bonferroni
adjustment across all four prespecified adjacent pairs gives
$p=1.14\times10^{-9}$. This supports the robustness of the reversal under a
paired-task inference model; it does not replace the original decision rule
or imply a random sample of deployed programming tasks.

\section{Mechanism Audit}
The retained StarCoder2 suffix for HumanEval/0 began by repeating the complete
function signature and docstring. EvalPlus registers the exact text
\verb|\ndef | as a generation stop for direct HumanEval completion, truncating
at that boundary and leaving an empty candidate body. Across 164 primary tasks,
142 retained suffixes were empty.

EvalPlus 0.3.1 constructs base prompts as
\verb|task["prompt"].strip() + "\n"|. We reran every model while retaining all
other inputs and stops but omitting that final newline. StarCoder2 rose from
3 to 49 HumanEval passes and from 3 to 42 HumanEval+ passes; 46 HumanEval tasks
changed fail-to-pass and none regressed. Every other checkpoint moved by at
most three net tasks. The published order was recovered ($\tau_b=1.0$,
task-resampling interval $[0.8,1.0]$). Because this condition was selected after
the primary result, it explains the local failure but does not replace the
primary decision or identify the unpublished historical source configuration.

We then completed the $2\times2$ newline-by-\verb|\ndef | stop factorial
(Table~\ref{tab:factorial}). Removing the newline eliminated empty suffixes.
Removing the stop exposed multi-definition continuation text in both prompt
conditions: with both the trailing newline and the exact stop removed, 143 suffixes contained a top-level
definition and 38 repeated the task entry point. Sanitization made the
functional effect far smaller. Of 164 no-newline/no-stop candidates, 151 were
byte-identical to the evaluated no-newline/standard-stop candidates; all 13
differences were baseline failures. Hardened evaluation scored the fourth cell
at 50/164 HumanEval and 43/164 HumanEval+, just one additional pass on each
suite. Thus the newline accounts for nearly all functional recovery, while the
stop mostly controls otherwise truncated continuation text.

\begin{table}[t]
\caption{StarCoder2 $2\times2$ factorial over 164 tasks. ``Def'' and ``repeat''
classify raw suffix text; HE is the number passing HumanEval.}
\label{tab:factorial}
\centering
\small
\begin{tabular}{lrrrr}
\toprule
Condition & Empty & Def & Repeat & HE \\
\midrule
Newline, stop & 142 & 0 & 0 & 3 \\
Newline, no stop & 2 & 152 & 63 & 17 \\
No newline, stop & 0 & 0 & 0 & 49 \\
No newline, no stop & 0 & 143 & 38 & 50 \\
\bottomrule
\end{tabular}
\end{table}

These definition counts are lexical: the classifier matches an unindented
\texttt{def} at the start of a line, and ``repeat'' matches the task entry-point
name. It is not an AST or execution-based classification and can match text
inside multiline strings. Zeros under the standard stop are post-truncation
observations, not evidence that the decoder never attempted another definition.
The newline/stop cell uses the optimized primary run (3 passes); the separate
unmodified-provider control has 2 passes and is not substituted into that cell.

\begin{table}[t]
\caption{All-model prompt-boundary check: primary versus no-newline passes.
Each denominator is 164; the condition is post hoc.}
\label{tab:newline}
\centering\small
\begin{tabular}{lrr}
\toprule
Model & HE (primary $\to$ altered) & HE+ \\
\midrule
Qwen-0.5B & $39\to37$ & $33\to31$ \\
StarCoder2-3B & $3\to49$ & $3\to42$ \\
DeepSeek-1.3B & $56\to54$ & $47\to46$ \\
Qwen-1.5B & $64\to67$ & $56\to56$ \\
Qwen-3B & $86\to85$ & $70\to71$ \\
\bottomrule
\end{tabular}
\end{table}

Table~\ref{tab:newline} reports all five counterfactual counts rather than only
the outlier. Net changes can mask offsetting transitions: for example,
Qwen-1.5B gains ten HumanEval tasks and loses seven. For StarCoder2, the
HumanEval+ change comprises 40 gains and one loss.

\subsection{Sampling sensitivity}
The retained secondary experiment draws 20 samples per task at temperature
0.2, top-p 0.95, seed 11, and the same 512-token cap. HumanEval sampled pass@1
is 37.6\% for Qwen-1.5B, 32.2\% for DeepSeek-1.3B, and 2.3\% for StarCoder2
under the stock prompt. A post-hoc StarCoder2 no-newline condition reaches
29.8\%; corresponding pass@5 values are 50.8\%, 41.3\%, 5.6\%, and 40.4\%.
The prompt effect therefore appears under both measured decoding regimes.
These are single-seed results, not estimates of variation across seeds.
The protocol's extra-seed trigger was not met: the closest adjacent sampled
pair differed by 5.37 points, above the five-point threshold, with no reversal
of the primary local order among those three models.

\section{Positioning and Reuse}
EvalPlus demonstrates that test adequacy changes functional-correctness
measurements~\cite{liu2023evalplus}; broader benchmarks target library use and
contamination-aware evaluation~\cite{zhuo2024bigcodebench,jain2024livecodebench}.
Our audit holds task tests fixed and varies the boundary between prompts and
executable candidates. The contribution is the measured size and ranking
consequence of this interaction, with a complete task-level counterfactual,
rather than a claim that prompt sensitivity itself is new.
Distinguishing repeated execution from reproducible conclusions follows
Bouthillier et al.~\cite{bouthillier2019reproducible}.
The upstream provider is itself versioned software~\cite{evalplus2026setup};
its historical prompt changes make the pipeline, including literal prompt
bytes, an essential unit of comparison.

The artifact supports three levels of reuse: recounting recorded evaluator
outcomes without executing generated code; regenerating rank analysis from
paired task outcomes; and regenerating/evaluating candidates using pinned
model revisions and the container launcher. The first two need no GPU.
The Dockerfile alone is insufficient isolation: the launcher applies the
network, mount, capability, memory, and process restrictions at runtime.

\balance
\section{Implications and Threats}
Benchmark values should be reported as properties of complete evaluation
systems. At minimum, reports should preserve exact prompt bytes, decoding,
stopping, sanitizer and evaluator revisions, checkpoint revisions, and
task-level outcomes. Prompt construction deserves tokenizer-boundary tests
across model families rather than visual inspection of source strings.

The study transports a historical 2024 comparison to one 2026 local stack. It
does not prove that the pinned revisions are the historical snapshots, identify
the source pipeline, estimate hardware variation, or generalize beyond
HumanEval's 164 Python tasks. The diagnostic is post hoc and no pre-truncation
logits were retained. Task-resampling intervals describe stability to task
composition, not uncertainty about this fully enumerated benchmark. Broader
benchmarks such as BigCodeBench \cite{zhuo2024bigcodebench} and LiveCodeBench
\cite{jain2024livecodebench} test different capabilities; our claim is limited
to rank transport under the disclosed HumanEval pipeline.

\section{Artifact and Conclusion}
An anonymized, history-free supplement contains model revisions, raw and
sanitized generations, hardened evaluator outputs and image digests,
task-level analyses, checksums, and reproduction instructions. Generated code
is untrusted and must be evaluated only in the documented container.

The published five-model order did not transfer to the specified local
pipeline. A one-byte prompt-boundary change raised StarCoder2 from 3/164 to
49/164 and restored the order; disabling the interacting stop added only one
more pass. Portable benchmark rankings therefore require portable evaluation
pipelines, not model names and aggregate percentages alone.

\paragraph*{AI-use disclosure.} A generative AI assistant assisted with drafting,
editing, code, and artifact checks. Automated checks reconcile the reported
counts with retained evaluator outputs. The author remains responsible for
the content, citations, and final submission.

\newpage
\bibliographystyle{IEEEtran}
\bibliography{../../references}
\end{document}
